\chapter{Conception}
\label{chap:conception}

\section{Architecture globale du système}

La \cref{fig:architecture} présente l'architecture générale de
\projectname. Le \textbf{Poller Service} interroge périodiquement le Wazuh
Indexer pour récupérer les nouvelles alertes~; chaque alerte est évaluée par
le \textbf{Policy Engine}, qui décide si elle mérite l'ouverture d'une
investigation. Le cas échéant, l'\textbf{Investigation Manager} crée
l'investigation (et associe/crée automatiquement l'actif concerné), puis la
confie à l'\textbf{AI Investigator}, qui exécute sa boucle agentique en
dialoguant avec un fournisseur de LLM et en appelant les outils nécessaires.
Toutes les actions, preuves et le verdict final sont enregistrés en base
MySQL. L'\textbf{API FastAPI} expose l'ensemble (investigations, politiques,
actifs, assistant conversationnel, authentification) au \textbf{tableau de
bord Vue.js}.

\begin{figure}[H]
\centering
\resizebox{\textwidth}{!}{%
\begin{tikzpicture}[
    box/.style={draw, rounded corners, fill=brandlight, minimum width=3.4cm,
                minimum height=1cm, align=center, font=\small},
    llmbox/.style={draw, rounded corners, fill=orange!15, minimum width=2.8cm,
                minimum height=1cm, align=center, font=\small},
    dbbox/.style={draw, rounded corners, fill=green!10, minimum width=2.8cm,
                minimum height=1cm, align=center, font=\small},
    arr/.style={-Stealth, thick},
    node distance=0.9cm and 1.1cm
]
\node[box] (indexer) {Wazuh Indexer\\ (OpenSearch)};
\node[box, below=of indexer] (poller) {Poller Service};
\node[box, below=of poller] (policy) {Policy Engine};
\node[box, below=of policy] (manager) {Investigation\\Manager};
\node[box, below=of manager] (investigator) {AI Investigator\\(boucle agentique)};
\node[llmbox, right=of investigator] (llm) {Fournisseur LLM\\(API ou local)};
\node[box, below=of investigator] (chatbot) {Assistant\\conversationnel};
\node[dbbox, left=2.8cm of investigator] (mysql) {Base MySQL};
\node[box, right=2.6cm of manager] (api) {API FastAPI};
\node[box, right=of api] (frontend) {Tableau de bord\\Vue.js};
\node[box, above=of api, fill=red!10] (auth) {Authentification\\(JWT)};

\draw[arr] (indexer) -- (poller);
\draw[arr] (poller) -- (policy);
\draw[arr] (policy) -- (manager);
\draw[arr] (manager) -- (investigator);
\draw[arr] (investigator) -- (llm);
\draw[arr] (llm) -- (investigator);
\draw[arr] (investigator) -- (chatbot);
\draw[arr] (chatbot) -- (llm.south);
\draw[arr] (investigator.west) -- (mysql.east);
\draw[arr] (manager.west) -| (mysql.north);
\draw[arr] (chatbot.west) -| (mysql.south);
\draw[arr] (api) edge[bend left=15] (frontend)
           (frontend) edge[bend left=15] (api);
\draw[arr] (mysql.north) |- ($(api.west)+(0,-0.35)$);
\draw[arr] (auth) -- (api);
\draw[arr] (indexer) -- (auth.north -| indexer);
\end{tikzpicture}}
\caption{Architecture globale de \projectname}
\label{fig:architecture}
\end{figure}

\section{Diagramme de cas d'utilisation}

Le \cref{fig:usecase} identifie les trois acteurs du système~: l'
\textbf{Analyste} (utilisateur humain authentifié), le \textbf{Poller}
(déclencheur automatique périodique), et l'\textbf{Agent IA} (exécuté en
tâche de fond pour le compte du système).

\begin{figure}[H]
\centering
\begin{tikzpicture}[
    actor/.style={draw, minimum width=1.6cm, font=\small, align=center},
    usecase/.style={draw, ellipse, fill=brandlight, minimum width=3.6cm,
                     minimum height=1cm, align=center, font=\small},
    arr/.style={-, thick}
]
\node[actor] (analyst) at (0,0) {\Large$\bigcirc$\\[-2pt]Analyste};
\node[actor] (poller) at (0,-9) {\Large$\bigcirc$\\[-2pt]Poller};
\node[actor] (ia) at (13,-4.5) {\Large$\bigcirc$\\[-2pt]Agent IA};

\node[usecase] (login)      at (4.5, 1.5)  {Se connecter\\(comptes Wazuh)};
\node[usecase] (dashboard)  at (4.5, 0)    {Consulter le\\tableau de bord};
\node[usecase] (detail)     at (4.5,-1.5)  {Consulter le détail\\d'une investigation};
\node[usecase] (manual)     at (4.5,-3)    {Déclencher une\\investigation manuelle};
\node[usecase] (chat)       at (4.5,-4.5)  {Discuter avec\\l'assistant IA};
\node[usecase] (policies)   at (4.5,-6)    {Gérer les politiques\\d'investigation};
\node[usecase] (assets)     at (4.5,-7.5)  {Gérer les actifs};
\node[usecase] (poll)       at (9, -9)     {Collecter les\\alertes Wazuh};
\node[usecase] (evaluate)   at (9, -10.5)  {Évaluer une alerte\\selon les politiques};
\node[usecase] (investigate) at (9, -4.5)  {Investiguer une\\alerte (boucle agentique)};
\node[usecase] (tools)      at (13.2,-2.5) {Appeler un outil\\(recherche, réputation IP...)};
\node[usecase] (verdict)    at (13.2,-6.5) {Produire un\\verdict argumenté};

\draw[arr] (analyst) -- (login);
\draw[arr] (analyst) -- (dashboard);
\draw[arr] (analyst) -- (detail);
\draw[arr] (analyst) -- (manual);
\draw[arr] (analyst) -- (chat);
\draw[arr] (analyst) -- (policies);
\draw[arr] (analyst) -- (assets);
\draw[arr] (poller) -- (poll);
\draw[arr] (poller) -- (evaluate);
\draw[arr] (evaluate) -- (manual);
\draw[arr] (ia) -- (investigate);
\draw[arr] (ia) -- (tools);
\draw[arr] (ia) -- (verdict);
\draw[arr] (ia) -- (chat);
\draw[arr, dashed] (manual) -- (investigate) node[midway, above, sloped, font=\tiny] {\guillemotleft include\guillemotright};
\end{tikzpicture}
\caption{Diagramme de cas d'utilisation}
\label{fig:usecase}
\end{figure}

\section{Diagramme de classes}

Le \cref{fig:classdiag} présente un extrait représentatif du modèle de
classes~: la hiérarchie des \texttt{Tool} (chaque outil que l'IA peut
appeler), la hiérarchie des \texttt{LLMProvider} (abstraction indépendante
du fournisseur de modèle), et leur utilisation par les deux points d'entrée
de raisonnement, \texttt{AIInvestigator} et \texttt{SecurityChatbot}.

\begin{figure}[H]
\centering
\resizebox{\textwidth}{!}{%
\begin{tikzpicture}

\begin{abstractclass}[text width=6.2cm]{Tool}{0,11}
  \operation{+ name() : str}
  \operation{+ description() : str}
  \operation{+ parameters() : dict}
  \operation{+ execute(**kwargs) : Any}
  \operation{+ coerce\_arguments(args : dict) : dict}
\end{abstractclass}

{\scriptsize
\begin{class}[text width=3.2cm]{SearchAlertsTool}{-9.6,4.5}
\end{class}
\inherit{Tool}

\begin{class}[text width=3.2cm]{CountAlertsTool}{-5.6,4.5}
\end{class}
\inherit{Tool}

\begin{class}[text width=3.2cm]{CheckIPReputationTool}{-1.6,4.5}
\end{class}
\inherit{Tool}

\begin{class}[text width=3.2cm]{SearchInvestigationsTool}{2.4,4.5}
\end{class}
\inherit{Tool}

\begin{class}[text width=3.2cm]{SearchAssetsTool}{6.4,4.5}
\end{class}
\inherit{Tool}
}

\begin{abstractclass}[text width=5.4cm]{LLMProvider}{15,11}
  \operation{+ chat(messages, tools, tool\_choice) : LLMResponse}
\end{abstractclass}

\begin{class}[text width=3.6cm]{APIProvider}{13,4.5}
\end{class}
\inherit{LLMProvider}

\begin{class}[text width=3.6cm]{LocalProvider}{17,4.5}
\end{class}
\inherit{LLMProvider}

\begin{class}[text width=5.6cm]{AIInvestigator}{2,-1}
  \operation{+ investigate(alert\_data) : dict}
  \operation{- \_run\_loop(...) : dict}
\end{class}

\begin{class}[text width=5.6cm]{SecurityChatbot}{13,-1}
  \operation{+ send\_message(session\_id, text) : dict}
  \operation{+ create\_session() : dict}
\end{class}

\unidirectionalAssociation{AIInvestigator}{uses}{1}{LLMProvider}
\unidirectionalAssociation{AIInvestigator}{uses}{*}{Tool}
\unidirectionalAssociation{SecurityChatbot}{uses}{1}{LLMProvider}
\unidirectionalAssociation{SecurityChatbot}{uses}{*}{Tool}

\end{tikzpicture}}
\caption{Diagramme de classes (extrait représentatif)}
\label{fig:classdiag}
\end{figure}

\section{Diagrammes de séquence}

\subsection{Investigation automatique d'une alerte}

Le \cref{fig:seq1} illustre le déroulement d'une investigation déclenchée
automatiquement par le \textit{Poller Service}, jusqu'à l'obtention d'un
verdict par la boucle agentique.

\begin{figure}[H]
\centering
\resizebox{\textwidth}{!}{%
\begin{sequencediagram}
  \newthread{poller}{:PollerService}
  \newinst[2]{policy}{:PolicyEngine}
  \newinst[2]{manager}{:InvestigationManager}
  \newinst[2]{investigator}{:AIInvestigator}
  \newinst[2]{llm}{:LLMProvider}
  \newinst[2]{tool}{:Tool}
  \newinst[2]{db}{:MySQL}

  \begin{call}{poller}{get\_alerts()}{policy}{alertes}
  \end{call}
  \begin{call}{poller}{evaluate(alerte)}{policy}{decision}
  \end{call}
  \begin{sdblock}{alt}{si decision = ouvrir investigation}
    \begin{call}{poller}{create\_investigation(alerte)}{manager}{investigation}
    \end{call}
    \begin{call}{manager}{investigate(investigation)}{investigator}{}
    \end{call}
    \begin{sdblock}{loop}{tant que verdict non atteint (max N étapes)}
      \begin{call}{investigator}{chat(messages, tools)}{llm}{tool\_calls}
      \end{call}
      \begin{call}{investigator}{execute(args)}{tool}{résultat}
      \end{call}
      \begin{call}{investigator}{store\_evidence(résultat)}{db}{}
      \end{call}
    \end{sdblock}
    \begin{call}{investigator}{chat(messages, tool\_choice=none)}{llm}{verdict}
    \end{call}
    \begin{call}{investigator}{store\_analysis(verdict)}{db}{}
    \end{call}
  \end{sdblock}
\end{sequencediagram}}
\caption{Séquence -- investigation automatique d'une alerte}
\label{fig:seq1}
\end{figure}

\subsection{Authentification d'un utilisateur}

Le \cref{fig:seq2} illustre le mécanisme d'authentification, fondé sur une
délégation de la vérification des identifiants au Wazuh Indexer plutôt que
sur une base de mots de passe propre à l'application (voir
\cref{sec:securite_conception}).

\begin{figure}[H]
\centering
\resizebox{0.95\textwidth}{!}{%
\begin{sequencediagram}
  \newthread{user}{:Analyste}
  \newinst[2]{front}{:Frontend (Vue)}
  \newinst[2]{api}{:API (/auth/login)}
  \newinst[2]{indexer}{:Wazuh Indexer}
  \newinst[2]{db}{:MySQL (users)}

  \begin{call}{user}{saisir identifiants}{front}{}
  \end{call}
  \begin{call}{front}{POST /auth/login}{api}{}
  \end{call}
  \begin{call}{api}{GET / (Basic Auth)}{indexer}{200 ou 401}
  \end{call}
  \begin{sdblock}{alt}{si 200 (identifiants valides)}
    \begin{call}{api}{SELECT/INSERT user}{db}{utilisateur}
    \end{call}
    \begin{callself}{api}{create\_access\_token()}{JWT}
    \end{callself}
    \begin{callself}{api}{réponse : JWT + profil}{}
    \end{callself}
  \end{sdblock}
  \begin{sdblock}{alt}{si 401 (identifiants invalides)}
    \begin{callself}{api}{réponse : 401 Unauthorized}{}
    \end{callself}
  \end{sdblock}
  \begin{call}{front}{requêtes suivantes (Bearer JWT)}{api}{}
  \end{call}
\end{sequencediagram}}
\caption{Séquence -- authentification d'un utilisateur}
\label{fig:seq2}
\end{figure}

\section{Modèle de données}

Le \cref{fig:erd} présente le modèle entité-association simplifié de la
base MySQL de la plateforme.

\begin{figure}[H]
\centering
\resizebox{\textwidth}{!}{%
\begin{tikzpicture}[
    entity/.style={draw, rounded corners, fill=brandlight, align=left,
                   font=\scriptsize, inner sep=6pt},
    rel/.style={-Stealth, thick},
    every node/.style={font=\scriptsize}
]
\node[entity] (users) at (0,8) {\textbf{users}\\id (PK)\\username\\email\\role\\is\_active};
\node[entity] (assets) at (7,8) {\textbf{assets}\\id (PK)\\hostname\\ip\_address\\wazuh\_agent\_id\\criticality};
\node[entity] (policies) at (14,8) {\textbf{investigation\_policies}\\id (PK)\\nom, condition, sévérité};

\node[entity] (inv) at (7,4) {\textbf{investigations}\\id (PK)\\alert\_id\\asset\_id (FK)\\policy\_id (FK)\\status, severity};

\node[entity] (evidence) at (0,0) {\textbf{investigation\_evidence}\\id (PK)\\investigation\_id (FK)\\source\_type, data};
\node[entity] (analysis) at (7,0) {\textbf{investigation\_analysis}\\id (PK)\\investigation\_id (FK)\\verdict, confidence};
\node[entity] (actions) at (14,0) {\textbf{investigation\_actions}\\id (PK)\\investigation\_id (FK)\\tool, reasoning, result};

\node[entity] (sessions) at (0,4) {\textbf{chat\_sessions}\\id (PK)\\title};
\node[entity] (messages) at (-6,0) {\textbf{chat\_messages}\\id (PK)\\session\_id (FK)\\role, content};

\draw[rel] (assets) -- node[above]{1} node[below]{1..*} (inv);
\draw[rel] (policies) -- node[above]{1} node[below]{0..*} (inv);
\draw[rel] (inv) -- node[left]{1} node[right]{0..*} (evidence);
\draw[rel] (inv) -- node[above]{1} node[below]{0..*} (analysis);
\draw[rel] (inv) -- node[left]{1} node[right]{0..*} (actions);
\draw[rel] (sessions) -- node[left]{1} node[right]{0..*} (messages);
\draw[rel, dashed] (users) -- node[above, sloped]{aucune table IA ne référence users} (0,4.9);
\end{tikzpicture}}
\caption{Modèle de données simplifié (MySQL)}
\label{fig:erd}
\end{figure}

\noindent\textit{Remarque de conception importante~:} la table
\texttt{users} n'est \textbf{jamais} interrogée par les outils exposés à
l'IA (ni pour l'agent d'investigation, ni pour l'assistant conversationnel).
Cette exclusion n'est pas un filtre applicatif contournable~: il n'existe
tout simplement \textbf{aucune classe \texttt{Tool}} permettant de lire
cette table (voir \cref{sec:securite_conception}).

\section{Choix technologiques}

Le \cref{tab:techstack} justifie les principaux choix technologiques.

\begin{table}[H]
\centering
\caption{Choix technologiques et justifications}
\label{tab:techstack}
\begin{tabular}{@{}p{3cm}p{3.6cm}p{7cm}@{}}
\toprule
\textbf{Couche} & \textbf{Technologie} & \textbf{Justification} \\
\midrule
Backend & Python, FastAPI & Typage, performance asynchrone, génération
automatique de documentation OpenAPI (\texttt{/docs}) \\
Accès données & SQLAlchemy / PyMySQL & Requêtes SQL explicites,
contrôle fin adapté à un schéma relatif ement simple \\
Base applicative & MySQL & Fiabilité, transactions ACID, disponibilité
large \\
Source d'alertes & Wazuh Indexer (OpenSearch), \texttt{opensearch-py} &
Réutilisation du moteur de recherche déjà déployé avec Wazuh \\
IA / LLM & Interface \texttt{LLMProvider} abstraite (implémentations
API compatible OpenAI et Ollama local) & Portabilité entre fournisseurs
distants (Groq, OpenAI...) et modèle auto-hébergé, sans changement de
code métier \\
Authentification & JWT + délégation au Wazuh Indexer & Évite la
duplication d'une base de mots de passe (\cref{sec:securite_conception}) \\
Frontend & Vue 3 (Composition API), Pinia, Vue Router, Tailwind CSS,
Vite & Écosystème léger, réactivité, rapidité de développement \\
\bottomrule
\end{tabular}
\end{table}

\section{Sécurité de la conception}
\label{sec:securite_conception}

Deux décisions de conception structurent la sécurité de la plateforme~:

\begin{enumerate}
  \item \textbf{Aucune base de mots de passe propre à l'application.}
        L'authentification (\cref{fig:seq2}) délègue la vérification des
        identifiants au Wazuh Indexer~: le mot de passe saisi n'est
        transmis qu'une seule fois, pour cette vérification, et n'est
        jamais stocké. Un jeton \textbf{JWT} signé, à durée de vie limitée,
        est ensuite émis par l'application pour authentifier les requêtes
        suivantes.
  \item \textbf{Exclusion structurelle des données d'identité pour l'IA.}
        Aucun outil (\texttt{Tool}) exposé à l'agent d'investigation ou à
        l'assistant conversationnel ne permet d'interroger la table
        \texttt{users}. Il ne s'agit pas d'un filtrage de requête, mais
        d'une absence totale de chemin d'accès~: quelle que soit la
        formulation de la demande d'un utilisateur au chatbot, aucun outil
        disponible ne peut retourner de données d'identité.
\end{enumerate}

Par ailleurs, la boucle agentique elle-même intègre plusieurs
\textit{garde-fous} destinés à borner son comportement (nombre maximal
d'étapes, taille maximale du contexte renvoyé au modèle, détection des
appels d'outils dupliqués, délais d'expiration réseau) — ces mécanismes sont
détaillés au \cref{chap:implementation}.
